
%--------------------------
\myframe{Asymptotic Behavior: preliminaries}{%

  \textit{In the following we are going to see some results related to the
  \textbf{predictor} dynamics. The filtered estimates are intimately related and
  their behavior is essentially the same.}

  \pause
  \medskip
  \textit{Also, for the sake of compactness, in the next slides on the
  asymptotic behavior of the predictor we will resort to a slightly modified
  notation using subscripts, i.e.,} $\xhat(t|t-1)=\xhat_t$, $P(t|t-1)=P_t$.

  \pause
  \medskip
  \textit{Finally, we assume $S=0$, i.e., uncorrelated noises, to simplify the
  equations, i.e. $F=A$, $\tilde{Q}=Q$, $G_t=K_t=AL_t$, $\Gamma_t=A-K_tC$.} 

  \pause
  \medskip
  \medskip
  We recall the predictor dynamics equal to
  %
  \begin{align*}
    \xhat_{t+1} &= A\xhat_t + K_t\left[\y_t - C\xhat_t\right]\\
    P_{t+1} &=  A\left[P_t - P_tC^T\left(CP_tC^T + R\right)^{-1}CP_t \right]A^T + Q\\
    K_t &= AP_tC^T\left[CP_tC^T + R\right]^{-1}
  \end{align*}

}


%--------------------------
\myframe{Asymptotic Behavior: error dynamics}{%

  Also, defining the estimation error
  \[
    \e_t = \x_t - \xhat_t.
  \]
  %
  \pause
  After the update, the error dynamics are governed by
  %
  \begin{align*}    
   \e_{t+1} &= (A-K_tC)\e_t + \v_t - K_t \w_t\\
            &= \Gamma_t\e_t + \text{noise terms}
  \end{align*}

  \pause
  \vspace{.1cm}
  \textbf{Essentially, the covariance $P_t$ and, consequently, the gain $K_t$
  and closed-loop error $\Gamma_t$ dynamics are governed by the resolution of
  the discrete-time Riccati Equation (R.E.)}
  $$
  P_{t+1} =  A\left[P_t - P_tC^T\left(CP_tC^T + R\right)^{-1}CP_t \right]A^T + Q
  $$

  \medskip
  \pause
  Therefore:
  \begin{center}
  \emph{asymptotic filter behavior $\Longleftrightarrow$ asymptotic Riccati behavior.}
  \end{center}

  \medskip
  \pause
  \textit{Before looking at the behavior for the general case of non-stationary
  processes, we consider the stationary case.}

}

%--------------------------
\myframe{Asymptotic Behavior: the stationary case}{

\textit{In case of stationary processes and, in particular, asymptotically
stable systems, i.e., characterized by asymptotically stable state-space matrix
$A$ it is possible to state a strong direct result, connecting KF and
Wiener-Kolmogorov theories as well as the asymptotic stability of the filter.}

\pause
\medskip

\begin{exampleblock}{}
  If $A$ is asymptotically stable then
  $$
    \xhat_t = \Elin[\x_t|\y_0,\ldots,\y_t] \underset{t_0\to-\infty}{\longrightarrow} \xhat_t^\infty
  $$
  where $\xhat_t^\infty$ is the Wiener-Kolmogorov predictor of $\x_t$ based on
  the infinite past history of $\y$.
  Then, there exists
  $
  \lim_{t_0\to -\infty} P_t = P_\infty
  $
  with $P_\infty$ solution of the Algebraic Riccati Equation (ARE).
  The prediction satisfies
  \begin{align*}
    \xhat_{t+1}^\infty &= \Gamma_\infty\xhat_t^\infty + K_\infty\y_t \\ 
    K_\infty &= AP_\infty C^T\left[ CP_\infty C^T + R \right]^{-1} \\
    \Gamma_\infty &= A - K_\infty C 
  \end{align*}  
  And if $(\Gamma_\infty, K_\infty)$ is reachable, then $\Gamma_\infty$ is
  asymptotically stable.
\end{exampleblock}

}

%--------------------------
\myframe{Asymptotic Behavior: structural assumptions}{

  In the more general case of non-stationary processes ($A$ not asymptotically
  stable), the behavior is studied through the behavior of $P_t$ (i.e. of the RE)

  \medskip
  Yet, a couple of assumptions are needed\dots

  \pause
  \begin{block}{Detectability}
    The pair $(A,C)$ is \emph{detectable}: i.e., every unstable mode of $A$ is
    observable.
  \end{block}
  %
  \pause
  \begin{block}{Stabilizability}
    The pair $(A, Q^{1/2})$ is \emph{stabilizable}: i.e., every unstable mode of
    $A$ is reachable through the process noise channel.
  \end{block}
  %

  \pause
  These conditions (checked via PBH tests\footnote{\url{https://en.wikipedia.org/wiki/Hautus_lemma}}) are natural because:
  %
  \begin{itemize}
    \item unobservable unstable modes cannot be corrected from measurements,
    \item unreachable unstable modes create pathological Riccati behavior.
  \end{itemize}

  \pause
  \begin{alertblock}{}
    $R\succ0$, detect. of $(A,C)$, stab. of $(A,Q^{1/2})$ $\Longrightarrow$ $\exists !$ stabilizing steady-state $P_\infty$
  \end{alertblock}

}

% %--------------------------
% \myframe{Asymptotic Behavior: Discrete algebraic R.E.}{
%   A steady-state covariance $P_\infty$ must satisfy the fixed-point equation
%   \[
%   P = A P A^\top + GQG^\top - A P C^\top (CPC^\top + R)^{-1} C P A^\top.
%   \]
%   This is the \alert{filtering DARE}. Equivalently, the steady-state gain is
%   \[
%   K_\infty = P_\infty C^\top (C P_\infty C^\top + R)^{-1}
%   \]
%   for the filtered form, or the corresponding predictor gain depending on notation.

%   Important facts:
%   \begin{itemize}
%     \item the DARE can have multiple symmetric solutions,
%     \item only one is typically \emph{stabilizing},
%     \item the stabilizing solution is the one relevant for the asymptotic Kalman filter.
%   \end{itemize}
%   The associated steady-state error dynamics become time-invariant.
% }

% %--------------------------
% \myframe{Asymptotic Behavior: main convergence result $P$}{
%   Assume:
%   \[
%   R \succ 0, \qquad (A,C) \text{ detectable}, \qquad (A,GQ^{1/2}) \text{ stabilizable}.
%   \]
%   Then the Riccati recursion converges, for any initial covariance $P_0 \succeq 0$, to the unique stabilizing solution $P_\infty$ of the DARE:
%   \[
%   P_k \to P_\infty.
%   \]
%   Consequently,
%   \[
%   K_k \to K_\infty.
%   \]
%   Moreover, the corresponding closed-loop estimation dynamics are asymptotically stable:
%   \[
%   \rho\bigl(A - K_\infty C\bigr) < 1
%   \]
%   (or the equivalent predictor-form matrix under the adopted indexing convention).

%   Interpretation:
%   \begin{itemize}
%     \item after transients, the filter ``forgets'' the initial covariance,
%     \item the gain becomes constant,
%     \item the estimation error reaches a stationary covariance regime.
%   \end{itemize}
% }

%--------------------------
\myframe{Asymptotic Behavior: fundamental theorem of KF theory}{

  \begin{exampleblock}{}
    Necessary and sufficient condition for:
    
    \begin{itemize}
      \item $\exists !$ a symmetric $P_\infty$ solution of the A.R.E.
      $$
      P = APA^T - APC^T(CPC^T + R)^{-1}CPA^T + Q
      $$
      \item the solution $P_\infty$ is stabilizing
      \item $\lim_{t\to\infty} P_t = P_\infty\quad \forall$ initial symmetric
      positive semidefinite $P_0$ 
    \end{itemize}

    is that $(A,C)$ being detectable and $(F,Q^{1/2})$ being stabilizable.
  \end{exampleblock}

  \pause
  \medskip
  Interpretation:
  \begin{itemize}
    \item after transients, the filter ``forgets'' the initial covariance,
    \item the gain becomes constant,
    \item the estimation error reaches a stationary covariance regime.
  \end{itemize}

}

%--------------------------
\myframe{Asymptotic Behavior: steady-state KF}{
  
  Analogously to the stationary
  case, once $P_k$ and $K_k$ have converged, one can replace the time-varying
  filter by the constant-gain filter
  \begin{align*}
    \xhat_{t+1|t} &= A \xhat_{t|t},\\
    \xhat _{t+1|t+1} &= \xhat_{t+1|t} + K_\infty (\y_{t+1} - C \xhat_{t+1|t}).
  \end{align*}

  \pause
  Equivalently, for the predictor
  \[
  \xhat_{t+1|t} = (A-K_\infty C)\xhat_{t|t-1} + K_\infty y_{t}.
  \]

  \pause
  Advantages of the constant-regime filter:
  \begin{itemize}
    \item lower online computational cost,
    \item simpler implementation in embedded systems,
    \item explicit LTI error dynamics,
    \item same asymptotic performance as the full Kalman filter.
  \end{itemize}

  Hence, for time-invariant models, the practical end-point is often a
  \emph{steady-state Kalman filter}.

}

%--------------------------
\myframe{Asymptotic Behavior: transient vs asymptotic phases}{
  Two phases are usually observed:
  
  \pause
  \begin{block}{Transient phase}
    $P_k$ depends on $P_0$ and the gain $K_k$ varies significantly. The filter is learning how much to trust model vs measurements.
  \end{block}

  \pause
  \begin{block}{Asymptotic phase}
    $P_k \approx P_\infty$ and $K_k \approx K_\infty$. The filter has entered a stationary regime.
  \end{block}
  
  \pause
  The tuning parameters shape this regime:
  \[
  Q \uparrow\ \Rightarrow\ K_\infty \uparrow \quad \text{(more trust in measurements)},
  \]
  \[
  R \uparrow\ \Rightarrow\ K_\infty \downarrow \quad \text{(more trust in the model)}.
  \]

  \pause
  The asymptotic gain is the precise mathematical embodiment of the
model/measurement tradeoff. }

%--------------------------
\myframe{Asymptotic Behavior: continuous-time counterpart}{
  For the continuous-time model
  \[
  \dot x = A x + G w, \qquad y = Cx + v,
  \]
  with white noises of intensities $Q$ and $R$, the covariance satisfies the differential Riccati equation
  \[
  \dot P = A P + P A^\top + GQG^\top - P C^\top R^{-1} C P.
  \]
  \pause
  At steady state, $P_\infty$ solves the continuous algebraic Riccati equation (CARE)
  \[
  A P + P A^\top + GQG^\top - P C^\top R^{-1} C P = 0.
  \]
  \pause
  Under the analogous detectability/stabilizability assumptions,
  \[
  P_t \to P_\infty,
  \qquad
  K_t=P_tC^\top R^{-1} \to K_\infty = P_\infty C^\top R^{-1}.
  \]
  \pause
  The constant-regime filter is then the steady-state Kalman--Bucy filter.
}

%--------------------------
\myframe{Asymptotic Behavior: take-away results}{
  \begin{enumerate}
    \pause
    \item The asymptotic behavior of the Kalman filter is governed by a R.E.
    \pause 
    \item Under
    \[
    (A,C) \text{ detectable}, \qquad (A,Q^{1/2}) \text{ stabilizable}, \qquad R\succ0,
    \]
    the Riccati recursion converges to a unique stabilizing solution $P_\infty$.
    \pause  
    \item The filter gain converges: $K_k \to K_\infty$
    \pause
    \item The estimation error covariance converges: $P_k \to P_\infty$
    \pause
    \item The asymptotic filter becomes a constant-gain linear filter with the same steady-state performance as the full time-varying Kalman filter.
  \end{enumerate}

  \vspace{.6cm}
  \texttt{4\_asymptotic\_behavior.ipynb}
}

